\chapter{Design of a Frequent Batch Auction Engine}
\label{ch:design}

This chapter answers RQ1. It decomposes the engine into five design
blocks, enumerates the options available within each block, and derives
from the theory of \Cref{ch:foundations} the strategic behaviour each
option should be expected to induce. \Cref{sec:tradeoff} consolidates
the result into a trade-off matrix and \Cref{sec:casestudies} tests the
taxonomy against systems that have actually been deployed.

\section{The Engine as a Set of Design Blocks}
\label{sec:blocks}

A matching engine is a deterministic function from a stream of
timestamped messages to a stream of trades and book states. A continuous
engine applies that function incrementally: every message is processed
to completion before the next is read. A batch engine interposes an
accumulation phase, so the function is applied once per interval to a
set of messages rather than once per message.

That single change forces five decisions that a continuous engine never
has to make, because in a continuous engine each of them is implied by
serial processing.

\begin{enumerate}
  \item \textbf{The batch interval $\tau$.} How long the accumulation
  phase lasts and whether its end is predictable.
  \item \textbf{Information disclosure.} What participants observe
  during accumulation.
  \item \textbf{Order flow segregation.} Whether all orders compete in
  one pool.
  \item \textbf{Matching, clearing and rationing.} How the clearing
  price is computed and how the long side is allocated.
  \item \textbf{Settlement and residual handling.} How the cleared batch
  is settled and what happens to what did not clear.
\end{enumerate}

\Cref{fig:pipeline} shows how the blocks compose into a single batch
cycle. The remainder of the chapter treats each block in turn.

\begin{figure}[htbp]
\centering
\fbox{\begin{minipage}{0.92\textwidth}
\vspace{0.4em}
\centering
\ttfamily\small
\begin{tabular}{c}
\textbf{[1] Accumulation phase, duration $\tau$}\\
orders, cancels and modifies buffered; nothing executes\\
$\downarrow$ \quad \emph{governed by Block 1 (interval) and Block 2 (disclosure)}\\[0.5em]
\textbf{[2] Gate closes}\\
buffer frozen; late messages routed to batch $n+1$\\
$\downarrow$\\[0.5em]
\textbf{[3] Partition into auctions}\\
one pool, or maker/taker and retail/institutional partitions\\
$\downarrow$ \quad \emph{governed by Block 3 (segregation)}\\[0.5em]
\textbf{[4] Clearing}\\
solve for $\clr$ maximising executed volume\\
$\downarrow$ \quad \emph{governed by Block 4 (clearing)}\\[0.5em]
\textbf{[5] Rationing}\\
allocate the long side among eligible orders\\
$\downarrow$ \quad \emph{governed by Block 4 (rationing)}\\[0.5em]
\textbf{[6] Settlement and residual handling}\\
transfers executed; unmatched volume rolled over or cancelled\\
$\downarrow$ \quad \emph{governed by Block 5}\\[0.5em]
\textbf{[7] Publication} $\rightarrow$ next batch begins
\end{tabular}
\vspace{0.4em}
\end{minipage}}
\caption[The batch cycle and its design blocks]{One batch cycle of a
frequent batch auction engine, annotated with the design block that
governs each stage.}
\label{fig:pipeline}
\end{figure}

\section{Block 1: The Batch Interval}
\label{sec:interval}

\subsection{The Length of \texorpdfstring{$\tau$}{tau}}

The interval length is the parameter that determines how much of the
speed advantage the mechanism neutralises. Its effect is
straightforward: a speed differential of $\delta$ between two
participants matters only if it changes which batch their orders land
in, which happens only if their submission decisions fall within
$\delta$ of a boundary. For a uniformly distributed submission time the
probability of this is $\delta/\tau$. With $\delta$ on the order of
microseconds, which is the scale at which the races documented by
\citet{aquilina2022quantifying} are decided, and $\tau$ on the order of
100~milliseconds, the probability is on the order of $10^{-4}$ to
$10^{-5}$. Speed has not been prohibited; it has been made almost
worthless.

The cost is delay. Every participant waits, in expectation, $\tau/2$
before its order can execute, and $\tau$ in the worst case. For an
investor with a horizon of minutes this is negligible. For a market
maker managing inventory against a moving reference price it is a real
exposure, because the price can move during the interval while the
submitted order cannot be withdrawn in time to avoid the resulting
clearing.

Three considerations bound $\tau$ from opposite directions.

\begin{itemize}
  \item \textbf{From below:} $\tau$ must exceed the dispersion of
  network latencies among participants by a comfortable margin.
  Otherwise, a participant with a shorter path is systematically more
  likely to make the current batch, and the race survives in attenuated
  form.
  \item \textbf{From above:} $\tau$ must remain short enough that
  information is not left unincorporated for long. \citet{zoican2021speed}
  show that the relationship between auction frequency and market
  quality is not monotone. Very long intervals accumulate stale orders
  and increase the risk that clearing occurs at a price the market has
  already left behind.
  \item \textbf{From the side:} $\tau$ interacts with the volume that
  accumulates per batch. A thin instrument in a short interval may
  accumulate no crossing orders at all, so that batches clear empty and
  the mechanism degenerates into a delayed continuous market. Interval
  length and instrument liquidity are therefore not independent
  choices.
\end{itemize}

\subsection{Fixed versus Randomised Closing}

If the closing time is known in advance, it becomes a coordination
point. A fast participant can hold its decision until the last feasible
moment, observe the maximum amount of information, and then submit or
cancel. This is boundary sniping, and it is the interval-scale
reappearance of the continuous-time race, with one difference: the
prize is not a single stale quote but participation in a batch whose
composition is now partly known.

\begin{itemize}
  \item \textbf{Fixed deterministic boundary.} Simple, auditable,
  trivially schedulable by participants, and compatible with
  wall-clock-aligned reporting. Maximally exposed to boundary sniping.
  \item \textbf{Randomised closing time.} The gate closes at a random
  time drawn from a distribution announced in advance, for example
  uniform on $[\tau_{\min}, \tau_{\max}]$. Waiting until the last
  moment now risks missing the batch entirely, which converts a
  dominant strategy into a gamble. The cost is that participants can no
  longer plan submission precisely, and that the randomisation source
  must be verifiable, which is a non-trivial requirement in a
  decentralised setting.
  \item \textbf{Volume- or event-triggered closing.} The auction closes
  when accumulated volume crosses a threshold, or on a market event.
  This adapts the interval to activity, which addresses the thin
  instrument problem, but makes the closing time endogenous to
  participant behaviour and therefore manipulable: a participant can
  trigger the auction by submitting volume.
  \item \textbf{Randomised extension.} A deterministic nominal close
  followed by a short random extension if late activity is detected.
  This is the design used in several equity closing auctions and is a
  compromise: predictable in the normal case, unpredictable exactly when
  predictability would be exploited.
\end{itemize}

Theory predicts a clear ordering with respect to boundary sniping and a
weaker, opposite ordering with respect to operational simplicity and
participant planning. The empirical counterpart, boundary concentration
of arrivals, is measurable in principle but not as a behavioural
response under a replay design, for the reason stated in
\Cref{tab:predictions}.

\section{Block 2: Information Disclosure}
\label{sec:disclosure}

The disclosure regime determines what participants can condition on when
they submit. \Cref{tab:disclosure} sets out the four regimes introduced
in \Cref{sec:transparency} together with the behaviour each induces.

\begin{table}[htbp]
\centering
\caption[Disclosure regimes and induced behaviour]{Disclosure regimes
during the accumulation phase and their predicted behavioural
consequences.}
\label{tab:disclosure}
\small
\begin{tabularx}{\textwidth}{@{}p{2.9cm} L L@{}}
\toprule
\textbf{Regime} & \textbf{What is published during accumulation} &
\textbf{Predicted behaviour}\\
\midrule
Fully open book & All accumulated orders, continuously updated &
Strong within-batch price discovery; early submitters expose intent and
are exploited by later ones; incentive to submit as late as possible,
which compounds boundary sniping\\
Indicative clearing price & The price at which the batch would clear if
it closed now & Coordination on a common price estimate without exposing
individual orders; some intent leakage through price movement; the
regime of most equity opening and closing auctions\\
Aggregate imbalance & Net excess demand or supply, no price schedule &
Signals direction of pressure and attracts offsetting liquidity without
revealing individual limits; weaker price anchoring\\
Sealed bid & Nothing until the auction clears & No exploitation of
submitted orders and no incentive to delay for informational reasons;
clearing price uncertainty induces bid shading, which can reduce cleared
volume\\
\bottomrule
\end{tabularx}
\end{table}

Two points deserve emphasis because they are easy to miss.

First, disclosure and interval design are not independent. An open book
during accumulation strengthens the incentive to submit late, because
information accrues over the interval and a late submitter conditions on
more of it. An open book therefore raises the value of a randomised
close, whereas a sealed-bid design removes most of the informational
motive for waiting and makes a fixed boundary more defensible.

Second, opacity does not eliminate information leakage; it relocates it.
Under sealed bidding, the clearing price and volume of batch $n$ are
published before batch $n+1$ closes, so participants learn about
aggregate interest with a one-batch lag. As $\tau$ shrinks, the sequence
of published clearing prices approaches a continuous price signal, and
the difference between a sealed-bid FBA and a transparent one narrows in
informational terms even though it remains sharp within any single
batch.

\section{Block 3: Order Flow Segregation}
\label{sec:segregation}

\subsection{The Unified Pool}

In the default design, all orders enter one auction. This maximises the
set of potential counterparties and is the simplest to reason about and
to regulate. Its cost follows directly from \citet{glosten1985bid}: a
liquidity provider that cannot distinguish informed from uninformed flow
must quote a price reflecting the average information content of the
pool, so uninformed participants cross-subsidise the losses that
informed participants inflict.

\subsection{Segregation by Participant Class}

Splitting flow by participant class, typically retail against
institutional, allows a liquidity provider to quote a tighter price
against the subset it believes to be uninformed. This is the economic
logic behind retail wholesaling arrangements in equity markets, and it
transfers cleanly to a batch setting: run one auction per segment, or
one auction with segment-conditional eligibility.

The design questions it raises are mostly about classification. The
segmentation is only as good as the label, labels can be gamed by
routing institutional flow through retail channels, and segregation
raises fairness and best-execution questions because participants in
different segments receive different prices for the same asset at the
same instant. The Level~4 data used in this thesis supports
participant-level classification and therefore makes segment-conditional
analysis feasible in principle, though classification into benign and
toxic flow would itself require a behavioural criterion rather than an
observable label.

\subsection{Intent-Based Segregation: The Dual Flow Batch Auction}

Jump Trading's Dual Flow Batch Auction (DFBA) partitions by
\emph{intent} rather than by participant identity. Every order is
labelled as maker flow, meaning liquidity supply, or taker flow, meaning
executable demand. Within each batch the engine then runs two
independent auctions:

\begin{itemize}
  \item maker buys against taker sells, and
  \item maker sells against taker buys,
\end{itemize}

both settled at a single uniform clearing price. Two structural
consequences follow. Market makers never trade against each other,
which removes the inter-maker crossing that in a unified pool produces
volume without providing liquidity to anyone who wanted it. And a fast
participant cannot use taker flow to pick off maker flow inside the
batch, because the batch clears simultaneously at one price and there is
no queue to win.

The design carries its own costs. Splitting the pool splits the
liquidity, so each of the two auctions clears against a thinner book
than a unified auction would. Intent labelling must be enforced rather
than declared, since a participant that can label itself a maker in one
batch and a taker in the next captures the benefits of both. And the
partition presumes that maker and taker are meaningful categories, which
in a mechanism that clears everything simultaneously is a stipulation
imposed by the engine rather than a fact about arrival timing.

\subsection{The Maker--Taker Distinction and Fee Design}

This last point generalises into a problem the FBA literature has not
resolved. In a continuous book, maker and taker are defined by arrival
timing relative to book state: the resting order is the maker, the
crossing order is the taker. That definition is unavailable in a
mechanism where all matched orders arrive during the same interval and
execute at the same instant at the same price. Yet fee schedules,
rebates, and a substantial part of exchange revenue depend on the
distinction.

Four ways to redefine it are available, none of them fully
satisfactory:

\begin{enumerate}
  \item \textbf{By declared intent}, as in DFBA. Operationally simple,
  but self-declared unless enforced by a separate eligibility regime.
  \item \textbf{By limit price aggressiveness.} Orders whose limit is
  far from the clearing price on the passive side are makers; orders
  that would have crossed a wide range of prices are takers. This is
  observable and hard to game without changing economic exposure, but it
  taxes participants for demanding certainty of execution.
  \item \textbf{By marginality.} The orders that set the clearing price
  are takers; inframarginal orders are makers. Theoretically appealing,
  because it is exactly the marginal order that determines the price,
  but it makes the fee a participant pays depend on the actions of
  everyone else in the batch, which is unpredictable ex ante.
  \item \textbf{By persistence across batches.} An order that has rested
  through several batches is a maker; a new order is a taker. This
  rewards standing liquidity, which is the economic function the rebate
  is supposed to compensate, and it is the only one of the four that
  measures something continuous rather than instantaneous.
\end{enumerate}

The choice is not merely accounting. A rebate schedule is an incentive,
and whichever definition the engine adopts becomes the behaviour it
subsidises.

\section{Block 4: Matching, Clearing and Rationing}
\label{sec:clearing}

\subsection{The Single-Asset Uniform Clearing Problem}

Let a batch contain buy orders indexed by $i \in \mathcal{B}$ with limit
prices $\pi_i$ and quantities $q_i$, and sell orders indexed by
$j \in \mathcal{S}$ with limit prices $\pi_j$ and quantities $q_j$.
Aggregate demand and supply at a candidate price $p$ are the step
functions

\begin{equation}
D(p) \;=\; \sum_{i \in \mathcal{B} : \pi_i \geq p} q_i,
\qquad
S(p) \;=\; \sum_{j \in \mathcal{S} : \pi_j \leq p} q_j .
\label{eq:demandsupply}
\end{equation}

Executable volume at $p$ is $V(p) = \min\{D(p), S(p)\}$, and the
clearing price is any maximiser

\begin{equation}
\clr \in \argmax_{p \in \mathcal{P}} \; V(p),
\label{eq:clearing}
\end{equation}

where $\mathcal{P}$ is the set of admissible prices on the tick grid.
Because $D$ is non-increasing and $S$ is non-decreasing, $V$ is
unimodal and the maximiser is either unique or an interval. Exchanges
resolve the interval case with a documented cascade, and the choice of
cascade is itself a design decision:

\begin{enumerate}
  \item choose the price in the maximising set that minimises the
  absolute imbalance $|D(p) - S(p)|$;
  \item if still tied, choose the price on the side of the residual
  imbalance, so that pressure is reflected in the price;
  \item if still tied, choose the price closest to a reference, such as
  the previous batch's clearing price or an external oracle price.
\end{enumerate}

\subsection{The Equivalent Optimisation Formulation}

The same solution can be stated as a linear program over executed
quantities, which is the formulation that generalises to multiple assets
and which makes the economic interpretation of the clearing price
explicit. Let $x_i \in [0, q_i]$ be the executed quantity of buy order
$i$ and $y_j \in [0, q_j]$ that of sell order $j$:

\begin{align}
\max_{x, y} \quad & \sum_{i \in \mathcal{B}} \pi_i x_i \;-\; \sum_{j \in \mathcal{S}} \pi_j y_j
\label{eq:lpobj}\\
\text{s.t.} \quad & \sum_{i \in \mathcal{B}} x_i \;=\; \sum_{j \in \mathcal{S}} y_j
&& \text{(flow conservation)} \label{eq:lpflow}\\
& 0 \leq x_i \leq q_i \quad \forall i, \qquad
  0 \leq y_j \leq q_j \quad \forall j.
&& \text{(order bounds)} \label{eq:lpbounds}
\end{align}

The objective in \eqref{eq:lpobj} is total reported gains from trade,
and maximising it is equivalent to maximising executed volume under
standard conditions. The economically important observation is that the
dual variable associated with the flow conservation constraint
\eqref{eq:lpflow} is precisely the uniform clearing price $\clr$: it is
the shadow price of one additional unit of the asset in the batch.
Uniform pricing is therefore not an additional rule bolted onto a volume
maximisation; it is the dual of the constraint that defines the auction.

Complementary slackness gives the fill logic directly. Buy orders with
$\pi_i > \clr$ are fully filled, sell orders with $\pi_j < \clr$ are
fully filled, and only orders exactly at $\clr$ can be partially filled.
The rationing rule of \Cref{sec:rationing} therefore applies to a set
that is typically small but never economically negligible, since it
contains exactly the marginal participants.

\subsection{The Multi-Asset Extension}

In a venue that clears several assets simultaneously, as decentralised
batch exchanges do, the problem generalises. Let $\mathcal{A}$ be the
set of assets and let $\rho_a > 0$ be a price of asset $a$ in a common
numeraire. Flow conservation is imposed per asset,

\begin{equation}
\sum_{k \in \mathcal{K}} z_k \, \nu_{k,a} \;=\; 0 \qquad \forall a \in \mathcal{A},
\label{eq:multiflow}
\end{equation}

where $z_k$ is the executed fraction of order $k$ and $\nu_{k,a}$ is the
signed quantity of asset $a$ that order $k$ exchanges. The exchange rate
between any two assets is then $\rho_a / \rho_b$, which is coherent by
construction: the price vector admits no cycle of exchanges that returns
a profit, so the clearing is arbitrage-free across pairs. This is the
property that \citet{walther2021multitoken} and \citet{ramsay2023speedex}
identify as the central advantage of multi-asset batch clearing, and it
has no analogue in a continuous market, where each pair trades in its
own book and cross-rate consistency must be enforced by arbitrageurs
after the fact --- which is to say, by paying them. The computational
cost is that the multi-asset problem with discrete order acceptance is
combinatorial rather than a simple line search, which is why deployed
systems delegate it to competing solvers
\citep{harrison2025hybrid,canidio2026fair}.

The engines implemented for this thesis operate per asset pair, so the
single-asset formulation of \eqref{eq:clearing} is what is computed;
the multi-asset case is developed here because it belongs to the design
space and because it is the form the mechanism takes in the deployed
systems reviewed in \Cref{sec:casestudies}.

\subsection{Rationing Rules}
\label{sec:rationing}

At $\clr$ the two sides are in general not exactly equal. The long side
must be rationed, and the rule chosen determines what participants do to
improve their allocation.

\begin{description}
  \item[Price-time priority within the batch.] Orders at $\clr$ are
  filled in submission order. Preserves continuity with continuous
  market conventions and is easy to explain. It also reintroduces a
  speed race at interval scale: if being early inside the batch improves
  expected fill, participants invest in being early, which is precisely
  the incentive the mechanism was built to remove. This rule undoes part
  of the benefit of batching.
  \item[Pro rata.] Each order at $\clr$ receives a fill proportional to
  its size. Strategically neutral with respect to timing, which
  preserves the core property. It creates a different distortion:
  because allocation is increasing in submitted size, participants
  submit more size than they intend to trade, and the equilibrium
  involves inflated orders that participants hope will be scaled down.
  This is order stuffing, and it degrades the informational content of
  displayed size \citep{haynes2020precedence}.
  \item[Size priority or maximal fill.] Larger orders are filled first,
  or the allocation maximises the number of completely filled orders.
  Favours block liquidity, which may be the intended policy, but
  amplifies the size-inflation incentive of pro rata into a discrete
  threshold effect.
  \item[Random allocation.] Fills are assigned by lottery among orders
  at $\clr$. Removes both the timing and the size channel, and is the
  only rule under which submitted size is an unbiased statement of
  desired quantity. The cost is execution uncertainty that participants
  cannot manage, which is a real cost for anyone hedging, and it
  requires a randomness source that participants can verify.
  \item[Hybrid rules.] Production systems mix these. The CME's
  K-algorithm family allocates part of the volume by time priority and
  the remainder pro rata, with configurable weights, which lets the
  venue trade off the two distortions rather than choosing one.
  A batch-native analogue would allocate a share $\kappa$ pro rata and
  $1 - \kappa$ randomly, tuning $\kappa$ to balance predictability
  against size inflation.
\end{description}

\Cref{tab:rationing} summarises. The general result is that no rule is
neutral: every rule makes allocation a function of some variable, and
whichever variable that is becomes the one participants compete on.
The design question is therefore not which rule is undistorted but which
distortion the venue prefers.

\begin{table}[htbp]
\centering
\caption[Rationing rules and induced behaviour]{Rationing rules for the
marginal side of a batch and the behaviour each induces.}
\label{tab:rationing}
\small
\begin{tabularx}{\textwidth}{@{}p{3.0cm} L p{3.0cm}@{}}
\toprule
\textbf{Rule} & \textbf{Induced behaviour} & \textbf{Empirical
signature}\\
\midrule
Price-time within batch & Compete to submit early; micro-scale speed
race returns & Boundary concentration shifts to the start of the
interval\\
Pro rata & Inflate submitted size to increase share & Rising
submitted-to-filled size ratio\\
Size priority / maximal fill & Inflate size discretely to cross
thresholds & Clustering of order sizes just above thresholds\\
Random & No timing or size channel; execution uncertainty & Fill
distribution independent of size and arrival time\\
Hybrid ($\kappa$ pro rata, rest time or random) & Both channels present
but attenuated & Intermediate on both signatures\\
\bottomrule
\end{tabularx}
\end{table}

\section{Block 5: Settlement and Residual Handling}
\label{sec:settlement}

\subsection{Gross versus Net Settlement}

Once the batch has cleared, the transfers must be executed. Gross
settlement moves each matched pair separately, which is simple and
gives an unambiguous audit trail per fill. Net settlement computes each
participant's net position change and moves only that, which reduces the
number of non-zero transfers, sometimes dramatically when many
participants both buy and sell within the batch. In a decentralised
setting, where each transfer costs gas, the reduction is a direct cost
saving and is one of the main practical arguments for batching
\citep{wang2018loopring,gong2023fbadex}. In a centralised venue the
saving is smaller but the reconciliation burden is still lower. The
trade-off is that netting obscures the correspondence between individual
orders and individual transfers, which complicates per-fill reporting
and any fee schedule levied per execution.

\subsection{The Residual}

The volume on the heavier side that cannot be matched at $\clr$ does not
disappear and is not filled from anywhere else. The design question is
what happens to it, and there are exactly three answers, selected by the
order's time-in-force:

\begin{enumerate}
  \item \textbf{Roll over.} The unmatched remainder is carried into the
  next batch with its original limit price and, where relevant, its
  original submission timestamp. This is the natural default for
  resting limit orders and is the behaviour implemented in this thesis.
  \item \textbf{Cancel.} The remainder is discarded at the end of the
  batch, which is the batch analogue of an immediate-or-cancel order.
  \item \textbf{Cancel the whole order if not fully filled}, the analogue
  of fill-or-kill, which requires the constraint to be imposed inside
  the clearing problem rather than applied afterwards, because an
  all-or-nothing order makes \eqref{eq:lpobj} an integer program.
\end{enumerate}

The third option is worth pausing on, because it shows where the clean
linear structure of the clearing problem breaks. All-or-nothing
constraints turn a problem solvable by a line search over the price grid
into a combinatorial one, which is why venues that support them impose
size limits, and why decentralised batch exchanges that support rich
order types run solver competitions rather than a single deterministic
algorithm \citep{harrison2025hybrid}.

Crucially, no external liquidity source is used to absorb the residual.
The engines compared in \Cref{ch:simulation} can only execute against
order flow that was actually submitted, which is what makes the
comparison like-for-like: neither engine has access to liquidity the
other lacks.

\section{The Design Trade-off Matrix}
\label{sec:tradeoff}

\Cref{tab:tradeoff} consolidates the chapter. Each row is an option,
each option is scored on the four properties that the theory of
\Cref{ch:foundations} identifies as the relevant dimensions, and the
final column names the behaviour the option induces. The scoring is
ordinal and theoretical, not measured; it summarises the arguments made
above rather than adding to them.

\begin{sidewaystable}
\centering
\caption[The FBA design trade-off matrix]{The FBA design trade-off
matrix. Ratings are ordinal and derived from theory:
$+$ favourable, $\circ$ neutral or conditional, $-$ unfavourable.}
\label{tab:tradeoff}
\small
\begin{tabularx}{\textwidth}{@{}p{2.2cm} p{3.4cm} c c c c L@{}}
\toprule
\textbf{Block} & \textbf{Option} & \textbf{Speed neutral.} &
\textbf{Price disc.} & \textbf{Exec.\ certainty} & \textbf{Impl.\
simplicity} & \textbf{Induced behaviour}\\
\midrule
\multirow{4}{2.2cm}{Interval}
 & Short $\tau$ ($\leq 10$\,ms)      & $\circ$ & $+$    & $+$    & $\circ$ & Thin batches; residual speed value\\
 & Long $\tau$ ($\geq 500$\,ms)      & $+$     & $-$    & $-$    & $+$     & Stale orders; delay cost\\
 & Fixed boundary                    & $\circ$ & $+$    & $+$    & $+$     & Boundary sniping\\
 & Randomised close                  & $+$     & $\circ$& $-$    & $-$     & No timing target; planning cost\\
\midrule
\multirow{4}{2.2cm}{Disclosure}
 & Fully open book                   & $-$     & $+$    & $+$    & $+$     & Intent exposure; late submission\\
 & Indicative price                  & $\circ$ & $+$    & $\circ$& $\circ$ & Coordination without order-level leakage\\
 & Aggregate imbalance               & $+$     & $\circ$& $\circ$& $\circ$ & Directional signal only\\
 & Sealed bid                        & $+$     & $-$    & $-$    & $+$     & Bid shading\\
\midrule
\multirow{3}{2.2cm}{Segregation}
 & Unified pool                      & $\circ$ & $+$    & $+$    & $+$     & Cross-subsidy of informed flow\\
 & Class segregation                 & $\circ$ & $\circ$& $\circ$& $-$     & Tighter quotes to benign flow; label gaming\\
 & Intent segregation (DFBA)         & $+$     & $\circ$& $-$    & $-$     & No maker-vs-maker trades; split liquidity\\
\midrule
\multirow{5}{2.2cm}{Rationing}
 & Price-time in batch               & $-$     & $\circ$& $+$    & $+$     & Micro speed race\\
 & Pro rata                          & $+$     & $\circ$& $\circ$& $+$     & Order size inflation\\
 & Size priority                     & $+$     & $\circ$& $\circ$& $+$     & Threshold clustering\\
 & Random                            & $+$     & $\circ$& $-$    & $\circ$ & Unbiased size; execution uncertainty\\
 & Hybrid $\kappa$                   & $\circ$ & $\circ$& $\circ$& $-$     & Attenuated distortions\\
\midrule
\multirow{3}{2.2cm}{Settlement}
 & Gross transfers                   & $\circ$ & $\circ$& $+$    & $+$     & Full per-fill auditability\\
 & Net settlement                    & $\circ$ & $\circ$& $+$    & $-$     & Lower transfer cost; opaque per-fill mapping\\
 & Roll-over residual                & $\circ$ & $+$    & $\circ$& $+$     & Standing liquidity accumulates across batches\\
\bottomrule
\end{tabularx}
\end{sidewaystable}

\section{Case Studies: Deployed Batch Auction Systems}
\label{sec:casestudies}

The taxonomy is worth testing against systems that exist.
\Cref{tab:casestudies} places five of them in the design space.

\textbf{Equity opening and closing auctions} are the oldest deployment
and the most conservative point in the space: a long interval, an
indicative price, a unified pool, and price-time rationing, with a
randomised extension to protect the boundary. They demonstrate that
uniform-price clearing is operationally routine at scale, but their
frequency is once or twice per day, so they neutralise nothing about the
intraday speed race.

\textbf{The Taiwan Stock Exchange} operated intraday call auctions at
short intervals for decades before moving to continuous trading, which
makes it the closest thing to a natural experiment on frequent batching
in equities \citep{lee2026taiwan}.

\textbf{Jump Trading's DFBA} is the clearest instance of intent-based
segregation combined with batching, discussed in
\Cref{sec:segregation}.

\textbf{CoW Protocol} clears batches of user intents with a uniform
clearing price per token pair, delegating the clearing computation to a
competitive solver auction \citep{dutta2025cow,harrison2025hybrid}. It
is the deployed system in which the multi-asset formulation of
\eqref{eq:multiflow} does the most work, since coincidence-of-wants
matching within a batch avoids touching external liquidity at all.

\textbf{SPEEDEX} demonstrates that arbitrage-free multi-asset batch
clearing can be computed fast enough for a high-throughput venue, which
addresses the standard objection that batch clearing is too slow to be
practical \citep{ramsay2023speedex}.

\begin{table}[htbp]
\centering
\caption[Deployed batch auction systems in the design space]{Five
deployed systems placed in the design space of \Cref{sec:blocks}.}
\label{tab:casestudies}
\small
\begin{tabularx}{\textwidth}{@{}p{2.6cm} p{2.0cm} p{2.4cm} p{2.3cm} L@{}}
\toprule
\textbf{System} & \textbf{Interval} & \textbf{Disclosure} &
\textbf{Segregation} & \textbf{Clearing and rationing}\\
\midrule
Equity open/close auctions & Once or twice daily, randomised extension &
Indicative price and imbalance & Unified & Volume maximisation,
price-time rationing\\
TWSE intraday call auctions & Seconds to tens of seconds & Partial &
Unified & Volume maximisation\\
Jump DFBA & Sub-second batches & Not fully public & Maker/taker intent
partition & Two auctions per batch, one uniform price\\
CoW Protocol & Block-paced batches & Public intents & Unified, solver
mediated & Multi-asset uniform clearing via solver competition\\
SPEEDEX & Per ledger/block & Public & Unified & Multi-asset
arbitrage-free uniform clearing\\
\bottomrule
\end{tabularx}
\end{table}

\section{Answer to RQ1}
\label{sec:rq1answer}

RQ1 asked how a frequent batch auction engine is designed and which
options exist for each of its architectural components. The answer
developed in this chapter has three parts.

An FBA engine is fully specified by five design blocks: the batch
interval, the disclosure regime, the segregation architecture, the
clearing and rationing rule, and settlement and residual handling.
\Cref{fig:pipeline} shows that these blocks are sequential stages of one
batch cycle, and that the corresponding decisions in a continuous engine
are not absent but implied by serial processing, which is why they
become visible only once batching is introduced.

Within each block the options are enumerable and their behavioural
consequences derivable. \Cref{tab:tradeoff} states them. The recurring
structure is that each option shifts competition from one variable to
another rather than eliminating it: randomising the close removes the
timing target but adds planning cost, opacity removes intent exposure
but adds bid shading, pro rata removes the timing channel but adds size
inflation. The only design choice in this space that is close to a free
improvement is uniform pricing itself, and even that has a cost, namely
that the marginal participants at $\clr$ must be rationed by some rule
that will in turn be gamed.

Finally, two questions in this space have no settled answer and are
contributions of this chapter rather than summaries of the literature.
The maker--taker distinction has no natural definition under
simultaneous uniform-price clearing, and four candidate redefinitions
were set out in \Cref{sec:segregation}, each of which subsidises
different behaviour. And the price improvement that uniform pricing
produces relative to submitted limits is not a fixed property of the
mechanism but a quantity whose size and distribution across participants
depends on the clearing and rationing rules chosen, which is precisely
what the empirical part of this thesis measures.
